\documentclass[12pt,a4paper]{article}
\usepackage[UTF8]{ctex}
\usepackage{geometry}
\usepackage{graphicx}
\usepackage{booktabs}
\usepackage{longtable}
\usepackage{array}
\usepackage{multirow}
\usepackage{amsmath}
\usepackage{hyperref}
\usepackage{xcolor}
\usepackage{tikz}
\usepackage{enumitem}
\usepackage{fancyhdr}
\usepackage{titlesec}
\usepackage{tabularx}

\geometry{left=2.5cm, right=2.5cm, top=2.8cm, bottom=2.5cm}

\pagestyle{fancy}
\fancyhf{}
\fancyhead[L]{\small 软件项目开发与实践}
\fancyhead[R]{\small 需求分析文档}
\fancyfoot[C]{\thepage}
\renewcommand{\headrulewidth}{0.4pt}

\titleformat{\section}{\Large\bfseries\color{blue!70!black}}{第\,\thesection\,章}{0.8em}{}
\titleformat{\subsection}{\large\bfseries}{}{0em}{}

\begin{document}

\begin{titlepage}
\centering
\vspace*{3cm}
{\Huge\bfseries\color{blue!70!black} 星露谷战斗模块}\\[0.8cm]
{\LARGE\color{blue!50!black} 俯视角像素风动作角色扮演游戏}\\[0.5cm]
{\large\color{gray} --- 以《星露谷物语》矿洞战斗为灵感 ---}\\[3cm]
{\Large 项目需求分析文档}\\[2cm]
\renewcommand{\arraystretch}{1.8}
\begin{tabular}{rl}
\textbf{课程名称：} & 软件项目开发与实践 \\
\textbf{项目名称：} & 星露谷战斗模块 \\
\textbf{姓\quad 名：} & 谢宇鹏 \\
\textbf{学\quad 号：} & 5120243390 \\
\textbf{开发模式：} & B模式（AI智能协作开发） \\
\textbf{技术栈：} & Python + Pygame \\
\end{tabular}
\vfill
{\small 本文档为项目完整需求规格说明，可作为开发团队或AI编码工具的输入。}
\end{titlepage}

\tableofcontents
\newpage

\section{引言}
\subsection{项目背景}
《星露谷物语》矿洞战斗系统以俯视角像素风画面、多武器实时战斗、怪物追逐为核心玩法，深受玩家喜爱。本项目提取其战斗与探索的核心玩法，剥离农场经营部分，打造一款以战斗和探索为核心的独立小游戏。
\subsection{项目目标}
开发一款基于Python + Pygame的俯视角像素风动作RPG小游戏，玩家在多个地图中探索、与怪物实时战斗、切换武器与魔法、升级属性。
\subsection{文档用途}
本文档可供开发团队或AI编码工具作为输入，指导游戏的设计与实现。

\section{用户群体分析}
\subsection{目标用户画像}
\begin{table}[htbp]
\centering
\begin{tabularx}{\textwidth}{lX}
\toprule
\textbf{维度} & \textbf{描述} \\
\midrule
年龄段 & 10--25岁 \\
性别 & 不限 \\
游戏经验 & 具有基础操作经验 \\
兴趣偏好 & 像素风游戏、RPG类型、星露谷/塞尔达爱好者 \\
典型场景 & 课余休闲，单局5--30分钟 \\
\bottomrule
\end{tabularx}
\end{table}
\subsection{用户核心需求}
\begin{enumerate}[leftmargin=2em]
    \item \textbf{即时可玩}：启动后快速进入游戏，主菜单简洁
    \item \textbf{操作直觉}：方向键移动、空格攻击、功能键切换武器
    \item \textbf{渐进难度}：初期怪物弱小，后期增强
    \item \textbf{视觉反馈}：攻击命中、受伤、升级有明确反馈
    \item \textbf{进度保存}：支持手动存档
\end{enumerate}

\section{游戏功能需求}
\subsection{游戏模式}
单人离线，俯视角2D像素风，实时动作战斗（非回合制）。

\subsection{故事大纲}
玩家扮演冒险者，在被怪物侵占的像素世界中穿越多个区域清除怪物。每个区域有不同种类和难度的怪物，通过战斗获得经验值提升实力。失败条件为生命值归零。

\subsection{场景设计}

\subsubsection{主菜单}
深色背景+浮动粒子，标题浮动动画，菜单选项高亮脉冲，背景音乐循环播放。上下键选择，回车确认，ESC退出。

\subsubsection{游戏关卡（核心场景）}
\begin{itemize}[leftmargin=2em]
    \item 地图系统：俯视角瓦片地图，多图层叠加（碰撞层/草地层/装饰层/实体层）
    \item 视窗系统：屏幕小于完整地图，视窗跟随玩家移动
    \item 玩家角色：像素精灵，12组动画（4方向$\times$3状态）
    \item 敌人：多种怪物，AI驱动，寻路追踪
    \item 可破坏物：草地等装饰物可被攻击摧毁
    \item 障碍物：树木、石墙等不可通行区域
    \item 传送点：触碰后切换地图，带渐入渐出效果
    \item HUD：血条、能量条、武器图标、经验值
    \item 升级菜单：暂停后分配经验值
\end{itemize}

\subsubsection{死亡界面}
暗红色背景，显示经验，R重开/ESC退出，红色粒子特效。

\subsubsection{存档对话框}
启动时检测存档，弹出对话框选择继续或新游戏。

\subsection{角色与装备}

\subsubsection{玩家属性}
\begin{table}[htbp]
\centering
\begin{tabularx}{\textwidth}{lXX}
\toprule
\textbf{属性} & \textbf{初始值/上限} & \textbf{作用} \\
\midrule
生命值 & 200/600 & 归零时死亡 \\
能量值 & 100/300 & 魔法消耗，自动恢复 \\
攻击力 & 15/30 & 影响伤害 \\
魔法 & 5/15 & 影响魔法伤害和回复 \\
速度 & 300/720 & 移动速度 \\
\bottomrule
\end{tabularx}
\end{table}

\subsubsection{武器系统}
\begin{table}[htbp]
\centering
\begin{tabularx}{\textwidth}{lccX}
\toprule
\textbf{武器} & \textbf{伤害} & \textbf{冷却} & \textbf{定位} \\
\midrule
剑 & 15 & 短 & 均衡型 \\
长矛 & 30 & 长 & 高伤慢速 \\
斧头 & 20 & 中 & 过渡武器 \\
刺剑 & 8 & 极短 & 连击流 \\
叉 & 10 & 短 & 灵活型 \\
\bottomrule
\end{tabularx}
\end{table}

\subsubsection{魔法系统}
\begin{table}[htbp]
\centering
\begin{tabularx}{\textwidth}{lXX}
\toprule
\textbf{魔法} & \textbf{效果} & \textbf{说明} \\
\midrule
治疗 & 回复50HP & 消耗10能量，有粒子特效 \\
火焰 & 前方范围伤害 & 消耗20能量，5个火焰粒子 \\
\bottomrule
\end{tabularx}
\end{table}

\subsubsection{敌人系统}
\begin{table}[htbp]
\centering
\begin{tabularx}{\textwidth}{lccclX}
\toprule
\textbf{怪物} & \textbf{血量} & \textbf{伤害} & \textbf{经验} & \textbf{攻击} & \textbf{AI特点} \\
\midrule
竹子 & 35 & 5 & 120 & 叶子 & 入门级 \\
鱿鱼 & 50 & 15 & 100 & 斩击 & 中速 \\
精灵 & 40 & 8 & 110 & 雷击 & 快速低伤 \\
浣熊 & 120 & 25 & 250 & 爪击 & 血厚高伤 \\
\bottomrule
\end{tabularx}
\end{table}

\subsection{地图与场景切换}
四图层CSV地图，传送点触发切换，淡出→加载新地图→淡入，玩家状态跨地图保留。

\subsection{存档系统}
JSON格式，保存玩家位置/血量/能量/经验/属性/装备/已击败敌人/已破坏草地。启动时检测并弹出对话框。

\subsection{升级系统}
消耗经验值提升五项属性，每项有上限。

\section{非功能需求}
\begin{table}[htbp]
\centering
\begin{tabularx}{\textwidth}{lX}
\toprule
\textbf{指标} & \textbf{要求} \\
\midrule
帧率 & 60FPS以上 \\
分辨率 & 1280$\times$720 \\
启动时间 & 3秒内 \\
输入响应 & 不超过1帧 \\
语言 & Python 3.10+ \\
引擎 & Pygame 2.x \\
架构 & 面向对象，多文件 \\
\bottomrule
\end{tabularx}
\end{table}

\section{可扩展性设计}
\begin{enumerate}[leftmargin=2em]
    \item 新怪物：添加配置数据和动画资源即可
    \item 新武器：添加武器配置，美术资源按方向提供
    \item 新魔法：添加配置和释放逻辑
    \item 新地图：创建CSV文件，配置传送点映射
    \item 商店系统：利用经验/金币扩展经济系统
    \item Boss战：创建Boss子类，多阶段血条
    \item NPC对话：定义NPC类型，对话框选项分支
\end{enumerate}

\section{验收标准}
\begin{enumerate}[leftmargin=2em]
    \item 游戏正常启动，主菜单显示且音乐播放
    \item 玩家四方向移动，动画正确
    \item 五种武器和两种魔法可用
    \item 四种敌人AI行为正常（待机/追踪/攻击）
    \item 击杀获得经验，可升级属性
    \item 多地图切换，渐入渐出效果
    \item 存档/读档正常
    \item 死亡后可重新开始
    \item 帧率60FPS+
    \item 面向对象架构清晰
\end{enumerate}

\section*{参考文献}
\addcontentsline{toc}{section}{参考文献}
\begin{enumerate}[leftmargin=2em]
    \item Pygame Official Documentation. \url{https://www.pygame.org/docs/}
    \item Stardew Valley. ConcernedApe.
    \item The Legend of Zelda Series. Nintendo.
    \item A* Pathfinding Introduction. Red Blob Games.
    \item 软件项目开发与实践课程资料. 刘畅.
\end{enumerate}

\end{document}
